\clearpage
\begin{proofbox}
  \textbf{\textcolor{CProof}{要证明的结论}}

  已知 $f(x)$ 是奇函数，并且 $0$ 在定义域内，证明 $f(0)=0$。
  下面的三种讲法，本质上都利用了奇函数的“相反数”规则；
  第一种算得最快，第二种最直观，第三种说明为什么不可能取其他数。

  \medskip
  \textbf{\textcolor{CProof}{证明一：直接代入奇函数的规则（最快，推荐）}}
\begin{center}\begin{tikzpicture}[font=\small,>=stealth,line join=round]
\draw[gray!40,rounded corners=3pt] (0,0) rectangle (14.2,8.5);\node at (7.1,7.9) {代入零以后，函数值必须是自己的相反数};
\draw[->,gray] (1,5.5)--(13.1,5.5);\foreach \x/\v in {3/-2,7.1/0,11.2/2}{\draw[gray] (\x,5.5)--(\x,5.35);\node[below=6pt] at (\x,5.5){$\v$};}\fill[blue] (3,5.5) circle(2.3pt);\fill[orange!85!black] (11.2,5.5) circle(2.3pt);\draw[->,blue] (11.2,5.7) .. controls (10,6.8) and (4,6.8) .. (3,5.7);\node at (7.1,6.9){取相反数：$2\mapsto-2$};\node at (7.1,4.35){$2$ 与 $-2$ 是两个不同的数};\draw[->,gray] (1,2.5)--(13.1,2.5);\fill[orange!85!black] (7.1,2.5) circle(2.5pt);\node[below=5pt] at (7.1,2.5){$0$};\draw[->,blue] (7.4,2.6) .. controls (9,3.7) and (5.2,3.7) .. (6.8,2.6);\node at (7.1,1.4){只有 $0$ 取相反数后仍在原位置：$0=-0$};
\node[text width=12.8cm,align=center] at (7.1,.65) {两条数轴比例相同；$f(0)=-f(0)$ 正是下方零点的自反关系。};
\end{tikzpicture}\end{center}


  \begin{description}[style=nextline, leftmargin=2.8em, labelsep=0.5em, itemsep=0.55em]
    \item[\textbf{第一步：写出规则。}]
      因为 $f$ 是奇函数，所以定义域内的每一个 $x$ 都满足
      \[
        f(-x)=-f(x).
      \]
    \item[\textbf{第二步：让 $x=0$。}]
      题目已经保证 $0$ 可以代入，而且 $-0=0$，因此
      \[
        f(-0)=-f(0)\quad\Longrightarrow\quad f(0)=-f(0).
      \]
    \item[\textbf{第三步：解这个最简单的方程。}]
      两边加上 $f(0)$：
      \[
        2f(0)=0\quad\Longrightarrow\quad\boxed{f(0)=0}.
      \]
  \end{description}

  \medskip
  
\end{proofbox}
\clearpage
\begin{proofbox}
\textbf{\textcolor{CProof}{证明二：用图像和“一个 $x$ 只能对应一个 $y$”（最直观）}}
\begin{center}\begin{tikzpicture}[font=\small,>=stealth,line join=round]
\draw[gray!40,rounded corners=3pt] (0,0) rectangle (14.2,8.5);\node at (7.1,7.9) {在同一条竖线上，对称点必须重合};
\node at (3.5,6.85){假设 $c\ne0$：同一个输入出现两个值};\node at (10.6,6.85){$c=0$：两个点合为原点};\begin{scope}[shift={(3.5,3.8)},x=1cm,y=1cm]\draw[->,gray] (-2.1,0)--(2.1,0) node[right]{$x$};\draw[->,gray] (0,-2.0)--(0,2.0) node[above]{$y$};\fill[red] (0,1.3) circle(2.2pt) node[right=5pt] {$(0,c)$};\fill[red] (0,-1.3) circle(2.2pt) node[right=5pt] {$(0,-c)$};\draw[red,dashed] (0,-1.3)--(0,1.3);\end{scope}\begin{scope}[shift={(10.5,3.8)},x=1cm,y=1cm]\draw[->,gray] (-2.1,0)--(2.1,0) node[right]{$x$};\draw[->,gray] (0,-2)--(0,2) node[above]{$y$};\draw[blue,thick,domain=-1.2:1.2,samples=150] plot(\x,{(\x)^3});\fill[orange!85!black] (0,0) circle(2.2pt) node[below right=5pt] {$(0,0)$};\end{scope}\node at (3.5,1.7){左侧不能作为函数的图像};\node at (10.6,1.7){一个横坐标只对应一个纵坐标};
\node[text width=12.8cm,align=center] at (7.1,.65) {左图为 $c\ne0$ 的假设示意，并非函数图像；$c<0$ 时上下两点交换。};
\end{tikzpicture}\end{center}


  \begin{description}[style=nextline, leftmargin=2.8em, labelsep=0.5em, itemsep=0.55em]
    \item[\textbf{第一步：设图像在 $y$ 轴上有一点。}]
      因为 $f(0)$ 有定义，可把这个点写成 $(0,c)$，其中 $c=f(0)$。
    \item[\textbf{第二步：找到它关于原点的对称点。}]
      奇函数的图像关于原点对称，所以 $(0,-c)$ 也必须在图像上。
    \item[\textbf{第三步：利用函数的基本规则。}]
      $(0,c)$ 和 $(0,-c)$ 的横坐标都是 $0$，
      一个横坐标不能对应两个不同的函数值。因此
      \[
        c=-c\quad\Longrightarrow\quad c=0.
      \]
      故 $\boxed{f(0)=0}$。
  \end{description}

  \medskip
  
\end{proofbox}
\clearpage
\begin{proofbox}
\textbf{\textcolor{CProof}{证明三：先假设它不等于零，再看哪里出错（反过来想）}}
\begin{center}\begin{tikzpicture}[font=\small,>=stealth,line join=round]
\draw[gray!40,rounded corners=3pt] (0,0) rectangle (14.2,8.5);\node at (7.1,7.9) {非零截距会怎样破坏原点对称？};
\begin{scope}[shift={(4.0,3.8)},x=1.2cm,y=0.9cm]\draw[->,gray] (-1.9,0)--(1.9,0) node[right]{$x$};\draw[->,gray] (0,-2.8)--(0,2.8) node[above]{$y$};\draw[blue,thick,domain=-1.05:1.05,samples=150] plot(\x,{(\x)^3+1});\draw[violet,dashed,domain=-1.05:1.05,samples=100] plot(\x,{(\x)^3-1});\fill[blue] (0,1) circle(2.2pt) node[above left=5pt] {$(0,1)$};\draw[violet,fill=white] (0,-1) circle(2.5pt);\node[below right=8pt,violet] at (0,-1){$(0,-1)$};\end{scope}\node[align=left,anchor=west] at (8.0,4.1){蓝实线：$y=x^3+1$\\[11pt]紫虚线：旋转半圈后的图像\\[11pt]两图不重合，不能是奇函数\\[11pt]只看 $f(0)=1\ne0$ 就能排除};
\node[text width=12.8cm,align=center] at (7.1,.65) {虚线为蓝图的原点对称像 $y=x^3-1$，不是原函数的另一部分。};
\end{tikzpicture}\end{center}


  \begin{description}[style=nextline, leftmargin=2.8em, labelsep=0.5em, itemsep=0.55em]
    \item[\textbf{第一步：作一个相反的假设。}]
      假设 $f(0)=c$，并且 $c\ne0$。
    \item[\textbf{第二步：使用奇函数条件。}]
      因为 $0$ 在定义域内，奇函数的规则要求
      \[
        f(-0)=-f(0)=-c.
      \]
    \item[\textbf{第三步：发现矛盾。}]
      但是 $-0=0$，因此左边实际上就是 $f(0)=c$，
      于是 $c=-c$，得到 $c=0$，与一开始的 $c\ne0$ 矛盾。
      所以假设不成立，只能是 $\boxed{f(0)=0}$。
  \end{description}

  \medskip
  \textbf{\textcolor{CProof}{三种证明怎样选择？}}

  考试作答优先用\textbf{证明一}，两行就够；
  向初中生解释时优先用\textbf{证明二}；
  需要说明“为什么非零值一定不行”时用\textbf{证明三}。
  不论哪一种，\textbf{都要先确认 $0$ 在定义域内}。

  \medskip
  \textbf{\textcolor{CProof}{顺手得到的快速判断}}

  若 $0$ 有定义，而代入得出 $f(0)\ne0$，
  那么上面的推导在 $x=0$ 这一处已经无法满足奇函数条件，
  因此\textbf{可以直接判定函数不是奇函数}，不必检查其他位置。
\end{proofbox}
